\chapter{Enterprise Architectures \& Solution Design}

Traditional software struggles with unstructured inputs (e.g. emails, logs). Agents solve this by acting as translators between unstructured data and structured code executions. 

Below are three production-ready enterprise architectures, mapped with custom TikZ flowcharts.

\section{Use Case 1: Autonomous Customer Service Query Routing}

This architecture intercepts incoming customer support emails, extracts queries, executes sandboxed SQL queries to fetch order status, and generates personalized responses.

\begin{figure}[h]
\centering
\begin{tikzpicture}[
    scale=0.8, transform shape,
    node distance=1cm and 1.5cm,
    block/.style={draw, rectangle, rounded corners, minimum height=0.8cm, minimum width=3cm, align=center, fill=blue!10},
    decision/.style={draw, diamond, aspect=2.5, align=center, fill=yellow!10},
    arrow/.style={thick,->,>=stealth}
]

% Nodes
\node (ticket) [block] {User Support Ticket};
\node (router) [block, below=of ticket] {Router Agent (LLM)};
\node (intent) [decision, below=of router] {Intent Classification?};
\node (rag) [block, left=of intent, xshift=-1cm] {RAG Doc Retrieval};
\node (sql) [block, right=of intent, xshift=1cm, fill=red!10] {Sandboxed SQL Writer};
\node (run) [block, below=of sql] {Execute SQL Database Query};
\node (eval) [block, below=of run] {Result Evaluator Agent};
\node (comp) [block, below=of intent, yshift=-3.5cm, fill=green!10] {Response Compiler};

% Links
\draw[arrow] (ticket) -- (router);
\draw[arrow] (router) -- (intent);
\draw[arrow] (intent) -- node[above, font=\footnotesize] {Info Query} (rag);
\draw[arrow] (intent) -- node[above, font=\footnotesize] {Order Lookup} (sql);
\draw[arrow] (sql) -- (run);
\draw[arrow] (run) -- (eval);
\draw[arrow] (eval) -- (comp);
\draw[arrow] (rag) |- (comp);

\end{tikzpicture}
\caption{Autonomous Support Routing Architecture}
\end{figure}

\section{Use Case 2: DevSecOps Vulnerability Patching Pipeline}

This agentic pipeline connects to GitHub webhooks, parses static security alerts, uses an AST chunker to find buggy code, writes patches in an isolated sandbox, and runs tests before generating a PR.

\begin{figure}[h]
\centering
\begin{tikzpicture}[
    scale=0.8, transform shape,
    node distance=1cm and 2cm,
    block/.style={draw, rectangle, rounded corners, minimum height=0.8cm, minimum width=3.5cm, align=center, fill=green!10},
    arrow/.style={thick,->,>=stealth}
]

% Nodes
\node (webhook) [block, fill=gray!10] {GitHub Webhook (Alert)};
\node (ast) [block, below=of webhook] {AST Parser (Find bug scope)};
\node (agent) [block, below=of ast, fill=yellow!10] {Patching Agent (LLM)};
\node (sandbox) [block, below=of agent, fill=red!10] {Sandboxed Local Build};
\node (test) [block, below=of sandbox] {Unit Test Runner};
\node (gate) [block, below=of test, fill=blue!10] {Create PR (GitHub API)};

% Links
\draw[arrow] (webhook) -- (ast);
\draw[arrow] (ast) -- (agent);
\draw[arrow] (agent) -- (sandbox);
\draw[arrow] (sandbox) -- (test);
\draw[arrow] (test) -- (gate);

\end{tikzpicture}
\caption{DevSecOps Self-Healing Pipeline}
\end{figure}

\section{Use Case 3: Financial Market Aggregator}

This architecture compiles multi-page reports by dynamically fetching stock tickers, parsing quarterly PDF filings from SEC Edgar, executing code to generate charts, and outputting compiled reports.

\begin{figure}[h]
\centering
\begin{tikzpicture}[
    scale=0.8, transform shape,
    node distance=0.8cm and 2cm,
    block/.style={draw, rectangle, rounded corners, minimum height=0.7cm, minimum width=3.8cm, align=center, fill=purple!10},
    arrow/.style={thick,->,>=stealth}
]

% Nodes
\node (sec) [block] {SEC Edgar Edgar API};
\node (chart) [block, below=of sec] {Chart Generator (Sandboxed)};
\node (report) [block, below=of chart, fill=yellow!10] {Compilation Agent (LLM)};
\node (compile) [block, below=of report] {LaTeX pdflatex Compiler};
\node (output) [block, below=of compile, fill=green!10] {Output Market PDF};

% Links
\draw[arrow] (sec) -- (chart);
\draw[arrow] (chart) -- (report);
\draw[arrow] (report) -- (compile);
\draw[arrow] (compile) -- (output);

\end{tikzpicture}
\caption{Autonomous Financial Compilation Flow}
\end{figure}
